% !TeX program = pdflatex
\documentclass[11pt]{article}
\usepackage{lmodern}
\usepackage{booktabs}
\usepackage{amsmath}
\usepackage[margin=1in]{geometry}
\usepackage{graphicx}
\usepackage{float}
% Keep tables at their input position, beneath the relevant headings.
\floatplacement{table}{H}
\usepackage[justification=centering]{caption}
\usepackage[hidelinks]{hyperref}
\usepackage{tikz}
\usepackage{fvextra}

\title{How do sportsbooks differ in their pricing strategies across player-prop products} %and what role might behavioral biases play in these differences?
\author{Sam Officer}
\date{} % leave empty to hide date

\begin{document}
\raggedbottom

\section{Summary Statistics}

\subsection{Overall Summary by Bookmaker}

\input{../MetricOdds/overall_overround_by_bookmaker.tex}

\subsection{Product-Level Summaries by Bookmaker}

\input{../MetricOdds/latex_tables/overround_BetMGM.tex}

\input{../MetricOdds/latex_tables/overround_BetRivers.tex}

\input{../MetricOdds/latex_tables/overround_Caesars.tex}

\input{../MetricOdds/latex_tables/overround_DraftKings.tex}

\input{../MetricOdds/latex_tables/overround_FanDuel.tex}

\input{../MetricOdds/latex_tables/overround_William_Hill_US.tex}

% Start methodology after all summary tables.
\clearpage
\section{Methodology}

\subsection{Construction of the Quoted Overround}

The analysis measures combined Over/Under pricing using the excess
overround. Each observation pairs a sportsbook's Over and Under quotes
for the same game, player, prop market, threshold, and recorded observation
time. The baseline uses half-point thresholds to avoid integer-threshold
pushes. Incomplete pairs are excluded, while the original prices and
update timestamps are retained.

Let $A_{cbts}$ denote American odds for contract $c$, sportsbook $b$,
observation time $t$, and side $s \in \{O,U\}$. Decimal odds, including
the return of the original stake, are calculated as
\begin{equation}
d_{cbts} =
\begin{cases}
1 + \dfrac{A_{cbts}}{100}, & A_{cbts}>0, \\[6pt]
1 + \dfrac{100}{|A_{cbts}|}, & A_{cbts}<0.
\end{cases}
\end{equation}

Following the conventional calculation discussed by Hegarty and Whelan
\cite[p.~4, equation~(1)]{hegartywhelan2025}, the excess overround is
\begin{equation}
\label{eq:overround}
\operatorname{OR}_{cbt}
=
\frac{1}{d_{cbtO}}
+
\frac{1}{d_{cbtU}}
-1.
\end{equation}

For reporting and estimation, the variable is expressed in
percentage-point units:
\begin{equation}
\operatorname{OR}^{\mathrm{pp}}_{cbt}
=
100 \operatorname{OR}_{cbt}.
\end{equation}

For example, American odds of $-110$ on both sides produce decimal odds
of approximately $1.9091$, giving
$\operatorname{OR}_{cbt}=0.047619$ and
$\operatorname{OR}^{\mathrm{pp}}_{cbt}=4.7619$.
Hegarty and Whelan use the term overround for the sum of inverse decimal
odds; equation~\eqref{eq:overround} measures its excess above one.

The constructed variable summarizes the combined quoted pricing of the
two outcomes. Comparisons across sportsbooks match the game, player,
market, threshold, and recorded observation time, allowing differences
in the quoted overround to be examined for identical contracts.

\paragraph{Limitations.}
The excess overround is a quoted pricing measure rather than a direct
estimate of realized sportsbook profits or expected bettor losses.
Hegarty and Whelan \cite[pp.~5--7]{hegartywhelan2025} explain that the
conventional expected-return interpretation relies on equal expected
returns across outcomes and can fail under unequal margin allocation.
The measure also conceals how prices are distributed between Over and
Under, so the original side-specific prices are retained.
Exact-threshold matching restricts findings to shared contracts and may
exclude important threshold disagreements. Common observation timestamps
do not guarantee simultaneous updates or executable offers; quote ages
therefore require examination. Observed pricing differences alone do
not identify customer demand, risk exposure, or firms' motives.

% Keep the regression specification and its explanation together.
\clearpage
\subsection{Bookmaker and Product Differences in Overround}

To examine whether differences between DraftKings and FanDuel vary
across player-prop products, we estimate:
\begin{equation}
\label{eq:bookmaker_market}
\mathrm{Overround}_{i}
=
\alpha
+ \delta D_i
+ \sum_{m \in \mathcal{M}^{*}} \beta_m M_{im}
+ \sum_{m \in \mathcal{M}^{*}} \theta_m (D_i M_{im})
+ \varepsilon_i.
\end{equation}

Each observation $i$ is a complete Over/Under quote pair for a given
sportsbook, game, player, product, threshold, and recorded observation
time. The dependent variable is excess overround multiplied by 100,
so coefficient differences are measured in percentage points.
The sample contains observations from DraftKings and FanDuel.
The indicator $D_i$ equals one for DraftKings and zero for FanDuel.
The indicator $M_{im}$ equals one if observation $i$ belongs to
product $m$ and zero otherwise. Passing completions is the reference
product, and $\mathcal{M}^{*}$ contains the remaining products.

The intercept $\alpha$ is FanDuel's mean overround for passing
completions. The coefficient $\delta$ measures the DraftKings minus
FanDuel difference for completions. For each other product $m$,
$\beta_m$ measures its difference from completions at FanDuel.
The interaction coefficient $\theta_m$ measures how the
DraftKings--FanDuel difference for product $m$ differs from their
difference for completions. Consequently, the difference between
sportsbooks for product $m$ is $\delta+\theta_m$.

This specification describes differences in observed overrounds.
Without additional controls or contract matching, the estimates may
also reflect differences in the games, players, thresholds, and
observation times represented in each sportsbook's product sample.
They do not identify causal effects or the reasons for firms'
pricing decisions.

\begin{thebibliography}{9}

\bibitem{hegartywhelan2025}
Hegarty, T., and Whelan, K. (2025).
\textit{Estimating Expected Loss Rates in Betting Markets:
Theory and Evidence}.
Revised draft, April 2025, University College Dublin.
\url{https://www.karlwhelan.com/Papers/Overround.pdf}.

\end{thebibliography}

\end{document}
